% page 174 of the facsimile (image p-173 of the scan)
% block 157.5 x 246.3 mm ; left margin 8.0 mm
\pgc{-12.700mm}{0.000mm}{
\l{\cel{5}164.}
\l{}
\l{}
\l{\sou{fideikomiso}.\cel{2}(\sou{yuro-cienco}). Donaco o legaco ad ulu, konfidi-}
\l{\cel{3}ta a lu kom honesto, por transmisesor ad altru.-DFIS}
\l{}
\l{\sou{fidela}. Qua ne falias pri to quon expektas de lu persono a}
\l{\cel{3}qua lu ligesas per su-obligir. - EFIS.}
\l{}
\l{\sou{fiera}. Ta qua havas alta (ma generale justa) opiniono pri sua}
\l{\cel{3}valoro, merito, povo, e qua manifestas lu extere. -FIS.}
\l{}
\l{\sou{fifro}. (\sou{muziko)}. Mikra fluto sis-trua, di qua la sono esas tre}
\l{\cel{3}akuta, uzita olim kom akompananto di la tamburo dum mar-}
\l{\cel{3}chi armeala. - DEF.}
\l{}
\l{\sou{figo}. (\sou{bot}.)Frukto di qua la pulpo esas mola e sukroza, pro-}
\l{\cel{3}duktata da arboro de la familio \textquotedbl{}urticei\textquotedbl{}, di qui la flori}
\l{\cel{3}esas asemblita en recevuyo komuna. - L.\cel{1}\sou{ficus}.\cel{2}DEFIRS.}
\l{}
\l{\sou{figuro}. I. Formo videbla di korpo. - II. (\sou{retoriko}).Formo}
\l{\cel{3}di la parolo, qua modifikas la expreso di la pensado. -}
\l{\cel{3}III. (\sou{geom}.) Reprezenturo grafika di la formi quin geome-}
\l{\cel{3}trio konceptas en spaco.}
\l{}
\l{\sou{figuranto}. La homo qua, en teatro, engajesis por plear la ro-}
\l{\cel{3}lo di la protagonisti qui ne parolas od esas tre acesora.}
\l{\cel{3}- DEFS.}
\l{}
\l{\sou{fijo}. Karto sur qua onu skribas la tituli di la libri di bi-}
\l{\cel{3}blioteko, noti, dokumenti, quan onu klasifikas alfabetale}
\l{\cel{3}o numerale, en buxi, por rekursar a lu kande onu bezonas}
\l{\cel{3}ta o ca libro, dokumento, e c.}
\l{}
\l{\sou{fiko}. (\sou{patol.}) Surkreskajo karnoza, kun pedunklo streta, e}
\l{\cel{3}somito inflata figatre.}
\l{}
\l{\sou{fikario}. (\sou{bot.}) Planto herbacea di qua la flori esas flava,}
\l{\cel{3}di la genero \textquotedbl{}ranunkuli\textquotedbl{}. - L.\cel{1}\sou{ranunculus ficaria}.}
\l{}
\l{\sou{fiktiva}. Di qua la valoro, la realeso, esas nur konvencionala.}
\l{\cel{3}- DEFIRS.}
\l{}
\l{\sou{filo}. Brino tenua, longa, ek texo-materio, metalo.-FIS.}
\l{}
\l{\sou{filamento}. Elemento longa e tenua ye\cel{1}\sou{qua}\cel{1}konsistas ula tisui}
\l{\cel{3}vejetala, animalala o minerala. - DEFIS.}
\l{}
\l{\sou{filandro}. Filo lejera qua flugetas en aero. - F.}
\l{}
\l{\sou{filantropo}. La homo qua amas homaro. - DEFIRS.}
\l{}
\l{\sou{filaso}. Materio texebla (kanabo, lino, e c.) ne ja filigita.}
\l{}
\l{}
\l{\sou{filatelio}. Cienco, studio, sercho di la posto-marki ( timbri}
\l{\cel{3}- DEFIS.\cel{42}postala)}
}
